%% sections/discussion.tex — single-partition version
\section{Discussion}
\label{sec:discussion}

We extended the DyAM multimodal attention framework with two contributions, a competitive (OvO) attention
mechanism with fewer parameters and an NLP-based encoding of tabular clinical features, and evaluated them on
response (AUC) and progression-free survival in a 247-patient cohort using one 10-fold cross-validation partition.
The experiments support four conclusions.

\paragraph{Fusing sources is what moves performance.}
The best single sources reached AUCs of $0.708$ (PD-L1~TPS), $0.701$ (CT radiomics), $0.676$ (genomics) and $0.615$ (TMB).
Fusing four sources raised the AUC to $0.798$ under the cooperative attention of the original framework
($+0.090$ over PD-L1~TPS), a value that this evaluation reproduced and that matches the $0.80$ of the original
publication \cite{Vanguri2022}, and the best model reached $0.824$ (95\% CI $0.765$--$0.883$), $+0.116$ above
PD-L1~TPS, the biomarker in routine use, and above the modest discrimination reported for PD-L1 and TMB
\cite{Lu2019,Rizvi2015}.
The same pattern held for survival: the association of the risk score with progression-free survival rose
from a C-index of $0.56$--$0.59$ for a single source to $0.59$--$0.66$ for six or seven sources, and the log-rank
$\chi^2$ from $1.9$--$16.8$ to $12.1$--$33.0$. Two endpoints therefore point to the same
conclusion, that the sources carry complementary information that attention-based fusion can combine.

\paragraph{The NLP encoding of clinical variables is better than numeric labs.}
This was the most consistent finding of the study. Replacing the 13 numeric laboratory and clinical values by their
sentence embedding raised the AUC in all six comparisons that were made: $+0.018$ and $+0.019$ in the two pathology arms
under cooperative attention, $+0.019$ and $+0.022$ in the two seven-source configurations under cooperative attention, and
$+0.033$ and $+0.030$ in the same two configurations under OvO attention.
That a model that sees the encoded clinical profile in place of the numeric vector gains in every comparison, with both
attention mechanisms and both pathology encodings, makes a chance explanation less likely than a systematic benefit.
A plausible reason is that a sentence embedding represents combinations of variables, for example advanced age together
with low albumin and poor performance status, that are not explicit in a 13-dimensional numeric vector and that the
attention module would otherwise have to learn from a small training set \cite{Reimers2019,Wang2020MiniLM}; our data do
not test this mechanism.
The AUC with the encoding was also above that of the same sources without any clinical variable
($+0.020$ for DyAM and $+0.024$ for OvO at the primary benchmark), and the general-purpose sentence encoder was better than
the domain-specific BioClinBERT (Supplementary Note~S2), so no clinical fine-tuning was needed.

\paragraph{OvO attention gives the best fusion with a quarter of the parameters.}
At the primary benchmark OvO improved on cooperative attention by $+0.017$ without the NLP encoding ($0.801$ against $0.783$)
and by $+0.021$ with it ($0.824$ against $0.803$), and in the two seven-source configurations with labs OvO was above
the cooperative model as well ($0.782$ against $0.772$ and $0.824$ against $0.803$ with NLP).
OvO with the NLP encoding gave the highest AUC of all 21 configurations while training $1{,}242$ parameters against
$4{,}969$ for the cooperative model. OvO therefore offers the best accuracy of the two mechanisms with a simpler model;
the parameter count is a possible explanation, which we did not test by shrinking the cooperative attention matrix.
For survival, OvO with the NLP encoding gave the strongest separation of risk groups at BM2 (hazard ratio $6.64$
against $4.63$ for DyAM, log-rank $\chi^2 = 23.9$ against $17.6$, 12-month time-dependent AUC $0.733$ against $0.698$),
and the hazard ratio was higher for OvO at BM1 ($5.83$ against $5.27$), TMB+PDL1 ($3.00$ against $2.17$) and
Rad+IHC-A+Gen+TMB+PDL1 ($7.04$ against $4.88$). The C-index was close for the two mechanisms ($0.633$ against $0.640$ at BM2),
so the benefit of OvO lies in the separation of high- and low-risk groups rather than in the ranking of patients.
The cooperative model remained the stronger of the two in a few configurations, such as Rad+IHC-A+Gen+PDL1
($0.798$ against $0.768$), where it also had the strongest survival association of the study (C-index $0.655$, hazard ratio $8.59$),
so the choice between mechanisms may depend on the combination of sources.

\paragraph{The approach is practical.}
Encoding a patient profile takes less than one second on CPU, needs no labelled clinical text, and can be applied to any
structured electronic health record field that can be written as a sentence, so it can be transferred to other structured
clinical datasets without a fine-tuned clinical language model.

\paragraph{Clinical encodings and survival.}
The clinical encodings did not change the survival metrics consistently: hazard ratios were $5.0$--$6.3$, C-indices $0.62$--$0.64$ and
log-rank $\chi^2$ above $17$ for every variant. The NLP encoding matched numeric labs on these metrics, so its advantage in this study is on the
response endpoint, and we do not claim that it captures prognostic information.

\paragraph{Performance estimates.}
A harder evaluation gave lower values: training on the patients with the earliest CT accession numbers and testing on the later 124
(a proxy for time, since calendar dates were not available) gave AUCs of $0.703$ ($0.587$--$0.819$) for OvO+NLP and $0.700$ ($0.587$--$0.812$) for
DyAM+NLP, with $0.707$ and $0.678$ at seven sources. The cross-validated AUC of $0.824$ is therefore best read as the performance
under the evaluation protocol of the original framework rather than the expected performance in new patients. Because about 60 variants were compared on
the same 247 patients, the highest figures also include some selection, which we estimate at $0.014$--$0.031$ AUC.

\paragraph{Limitations.}
First, all results come from one 10-fold cross-validation partition of a single-institution retrospective cohort, which limits generalisability;
the confidence intervals reflect patient sampling but not the choice of partition.
Second, with $n = 247$ (62 responders) the 95\% interval of an AUC is about $\pm 0.07$, and the differences between models, which are
$0.02$ AUC or smaller, are smaller than this cohort can resolve (about 1{,}200--1{,}350 patients would be needed to detect $0.02$ AUC with 80\% power;
Supplementary Note~S3), so they are described by direction and magnitude.
Third, results depended on the numerical environment (the same code gave $0.798$ on Linux and $0.788$ on Windows at the primary benchmark), so all results
are from one pinned environment.
Fourth, the sentence encoder was used without fine-tuning on NSCLC clinical text.
Fifth, calendar dates were not available, so temporal validation was approximated by CT accession order and covered two configurations.
Finally, PD-L1 TPS was scored with the Sauter method, which may differ from the SP142 or 22C3 assays used elsewhere, so PD-L1-derived
features should be interpreted with this in mind.
